
% SGA 3，报告 VII_B；英文重建本第一轮的中文译本。
% 范围： 节 4.1.1 全文 依据文献局部页 74,
% 止于下节之前： 节 4.1.2.
% 合并阅读本页： 587. 重版印刷页： 576.
% 范围内原版边码： 542.
% 范围内编者注： 134.
% 依据： Polo--Gille Exp7B-13oct24.pdf,
% SHA-256 D303690E0195654BBE7C3DAA5D89EA23E24CD709A3575217234AC19AE8836694.
% 参照了 Jacob Reinhold 的 CC BY 4.0 Markdown，作为对照来源。
% 原英文 Loop 1 阶段说明：剩余域箭头图为第一轮直接从 PDF 截取的原图。（当前中文工程已采用原生图。）
\subsubsection*{4.1.1}
\phantomsection
\label{sga3:VIIB:4.1.1}
\footnote{\label{sga3:VIIB:nde:134}编者纠正了原文，后者给出的包含关系为
\[
  (\widehat X/\widehat S)^{(p)}
  \subset
  \widehat{X^{(p)}}/\widehat S
\]
而非反向包含。他们还指出，下文不会用到本段。}
设 \(S\) 是素域 \(\mathbb F_p\) 上的概形，\(\eta:X\to S\) 是 \(S\)-概形。用 \(\operatorname{fr}(S)\) 表示 \(S\) 的“绝对”Frobenius 态射：它在底拓扑空间上诱导恒等映射，并将结构层的每个截面 \(f\) 映为 \(f^p\)；参见\ VII\textsubscript{A}，\SGARefLink{sga3:VIIA:4.1}{4.1}。令
\[
  X^{(p)}=X\times_{\operatorname{fr}(S)}S
\]
为相应的 \(S\)-概形\textup{（VII\textsubscript{A}，同上）}；因此其到 \(S\) 的结构态射为 \(\operatorname{fr}(S)\circ\eta\)。
令 \(\widehat S\) 为这样的形式概形：其底拓扑空间是赋予离散拓扑的 \(S\) 的闭点集，在闭点 \(s\) 处的局部环为 \(\mathcal O_{S,s}\) 关于有限余长度理想所定义的线性拓扑的分离完备化
\[
  \mathcal O_{\widehat S,s}
  =
  \widehat{\mathcal O}_{S,s}
\]
\textup{（参见\ \SGARefLink{sga3:VIIB:1.2.6}{1.2.6}）}。因此其仿射代数为
\[
  k=\mathcal A(\widehat S)
  =
  \prod_{s\in S_{\mathrm{cl}}}
  \widehat{\mathcal O}_{S,s}.
\]
回顾\textup{（同上）}，\(\widehat X/\widehat S\) 表示由满足
\[
  [\kappa(x):\kappa(s)]<\infty,
\]
的点 \(x\in X\) 构成的形式 \(k\)-簇，其中 \(s\) 是 \(x\) 在 \(S\) 中的像。其在 \(x\) 处的局部环，是 \(\mathcal O_{X,x}\) 关于如下理想 \(I\) 定义的线性拓扑的分离完备化
\[
  \mathcal O_{\widehat X/\widehat S,x}
  =
  \widehat{\mathcal O}_{X,x}
\]
：\(\mathcal O_{X,x}/I\) 作为 \(\mathcal O_{X,x}\)-模具有有限长度，从而作为 \(\mathcal O_{S,s}\)-模也具有有限长度。因此可以通过基变换构造形式 \(k\)-簇
\[
  (\widehat X/\widehat S)^{(p)}
  =
  (\widehat X/\widehat S)
  \widehat\otimes_{\pi}k.
\]
还可以考察形式 \(k\)-簇 \(\widehat{X^{(p)}}/\widehat S\)。其底集合由满足如下条件的点 \(x\in X^{(p)}=X\) 构成：若 \(s\) 表示 \(x\) 在 \(S\) 中的像，则复合
\begin{center}
  \phantomsection
  \label{sga3:VIIB:4.1.1:residue-field-frobenius-diagram}
  % SGA3-VII-DIAG-121
  % SGA3-VII-DIAG-121
% Authority: Exp7B-13oct24.pdf, local p. 74.
% Crop: figures/Exp7B_localp074_residue_field_frobenius_diagram.png
% Locator: tex/components/65_expose_VIIB_section411_formal_completion_frobenius_en.tex, Loop 1 line 89.
\(
\kappa(s)\xrightarrow{\ \operatorname{fr}\ }\kappa(s)
\xrightarrow{\ \eta^\natural\ }\kappa(x)
\)

\end{center}
使 \(\kappa(x)\) 成为 \(\kappa(s)\) 的有限扩张。这时对
\[
  \eta^\natural:\kappa(s)\longrightarrow\kappa(x),
\]
也有相同结论，所以 \(x\) 是 \(\widehat X/\widehat S\) 的点，且
\[
  \mathcal O_{\widehat{X^{(p)}}/\widehat S,x}
  =
  \widehat{\mathcal O}_{X,x}
  \widehat\otimes_{\pi}
  \widehat{\mathcal O}_{S,s}
  =
  \widehat{\mathcal O}_{X^{(p)},x}.
\]
第二个等号成立，是因为构造 \(X\mapsto\widehat X/\widehat S\) 与积交换\textup{（参见\ \SGARefLink{sga3:VIIB:1.2.6}{1.2.6}）}。
由此可得，作为形式子簇有
\[
  \widehat{X^{(p)}}/\widehat S
  \ \subset\
  (\widehat X/\widehat S)^{(p)}
\]
当且仅当对 \(S\) 的每个闭点 \(s\)，域 \(\kappa(s)\) 在 \(\kappa(s)^p\) 上的次数有限时，等号成立。例如，若 \(S\) 在满足
\[
  [\kappa:\kappa^p]<\infty.
\]
的域 \(\kappa\) 上局部有限型，就属于这种情形。
